\section{The cyclic control as an exact change of arithmetic interpolation volume}
\label{sec:cyclic-control-determinant-increments}

Retain a nonempty reflection-stable packet with every original order,
the tensor degree $k\geq1$, the cyclic algebra
$C=\mathbb C[S]/(\chi)$ of dimension $q\geq1$, and its fixed monic
remainder coordinates. Let $p_j(S)$ be the original monic orthogonal
polynomials on $S=k/2+iu$ for the measure $m_{h,k}(u)\,du$, and put
\[
 b_j=[p_j]_\chi,\qquad
 \omega_j=\int_{\mathbb R}|p_j(k/2+iu)|^2m_{h,k}(u)\,du>0,
 \quad K_N=\sum_{j=0}^N\frac{b_jb_j^*}{\omega_j},\quad G_N=K_N^{-1}.
 \tag{CV.1}
\]
Here $N\geq q-1$; the original $\omega_0=\mu_h^k$ remains.
All determinants below are taken in these same coordinates.
Set $K_{-1}=0$ and
\[
 \begin{gathered}
 D_j=\det K_j,\qquad
 \widetilde K_N=K_{N-1}+\frac{b_{N+1}b_{N+1}^*}{\omega_{N+1}},
 \quad\widetilde D_N=\det\widetilde K_N,\\
 a=b_N^*G_Nb_N,\quad d=b_{N+1}^*G_Nb_{N+1},\quad
 c=b_N^*G_Nb_{N+1}.
 \end{gathered}
 \tag{CV.2}
\]
The matrix $\widetilde K_N$ records the actual polynomial subspace
spanned by $p_0,\ldots,p_{N-1},p_{N+1}$, with its inherited norm:
the degree-$N$ orthogonal direction is omitted and the degree-$(N+1)$
direction is retained. It can be singular. No inverse for it is used.

\begin{proposition}[Exact adjacent-volume identities]
The following equalities hold, including $N=q-1$:
\[
 \begin{split}
 1-\frac{D_{N-1}}{D_N}&=\frac a{\omega_N},\\
 \frac{D_{N+1}}{D_N}-1&=\frac d{\omega_{N+1}},\\
 \frac{\widetilde D_N}{D_N}
 &=\left(1-\frac a{\omega_N}\right)
       \left(1+\frac d{\omega_{N+1}}\right)
      +\frac{|c|^2}{\omega_N\omega_{N+1}}.
 \end{split}
 \tag{CV.3}
\]
Consequently the original rank-two control allowance satisfies
\[
 \boxed{\quad
 (\epsilon_{h,k,N}^{\mathrm{cyc}})^2
 =\frac{\omega_{N+1}}{\omega_N}
 \frac{D_{N+1}+D_{N-1}-D_N-\widetilde D_N}{D_N}.
 \quad}
 \tag{CV.4}
\]
Equivalently, retaining the entire complex cross term,
\[
 \begin{split}
 (\epsilon_{h,k,N}^{\mathrm{cyc}})^2
 =\frac{\omega_{N+1}}{\omega_N}
 \left[
 \left(1-\frac{D_{N-1}}{D_N}\right)
 \left(\frac{D_{N+1}}{D_N}-1\right)
 -\frac{|c|^2}{\omega_N\omega_{N+1}}
 \right].
 \end{split}
 \tag{CV.5}
\]
\end{proposition}
\begin{proof}
The first $q$ monic polynomials give a triangular basis modulo $\chi$,
so $K_N$ is positive definite in the stated range. For an invertible
$K$ and columns $x,y$, direct block elimination of
$\left(\begin{smallmatrix}K&x\\-y^*&1\end{smallmatrix}\right)$
in its two possible orders gives
$\det(K+xy^*)=\det K(1+y^*K^{-1}x)$.
Applying this with $K=K_N$ and the displayed signed increments gives
the first two equations, even if the matrix after subtraction is singular.
For completeness the same block elimination with two columns proves
$\det(K+UV^*)=\det K\det(I_2+V^*K^{-1}U)$.
Take
\[
 U=\left(-\frac{b_N}{\sqrt{\omega_N}},
            \frac{b_{N+1}}{\sqrt{\omega_{N+1}}}\right),\qquad
 V=\left(\frac{b_N}{\sqrt{\omega_N}},
            \frac{b_{N+1}}{\sqrt{\omega_{N+1}}}\right).
\]
The positive square roots only specify this auxiliary factorization;
the original columns, coordinates and norms in (CV.1) are unchanged.
The resulting two-by-two determinant is
\[
 \det\begin{pmatrix}
 1-a/\omega_N&c/\sqrt{\omega_N\omega_{N+1}}\\
 -\bar c/\sqrt{\omega_N\omega_{N+1}}&1+d/\omega_{N+1}
 \end{pmatrix},
\]
which proves the third equation with its plus sign.
Subtract that equation from the sum of the first two determinant ratios
as ordered in (CV.4). The resulting value is
$(ad-|c|^2)/(\omega_N\omega_{N+1})$.
The original control formula is
$\epsilon^2=(ad-|c|^2)/\omega_N^2$, proving (CV.4)--(CV.5).
In particular its numerator is nonnegative by the Gram determinant
of the two original vectors $b_N,b_{N+1}$ in $G_N$.
For $N=q-1$, $K_{N-1}$ has rank $q-1$, so $D_{N-1}=0$ and
$a=\omega_N$; every equation still holds. For $q=1$ the two-vector
Gram determinant is zero, consistently with the trace-zero
one-dimensional control. There is no positive-dimensional determinant
or Gram inverse to apply when $h=1$; the coefficient quotient is then
zero and its unique maps have the split lift $\{\tau,e\}$.
\end{proof}

\subsection{The exact volume map and the reflected quartet bound}

Let $\mathcal P_N$ be the span of $p_0,\ldots,p_N$ with its original
integral norm, and let $T_N:\mathcal P_N\to C$ be reduction modulo
$\chi$. In the displayed polynomial basis its matrix is
$[b_0\ \cdots\ b_N]$, and its source Gram is
$\operatorname{diag}(\omega_0,\ldots,\omega_N)$.
The adjoint is therefore the explicit map
\[
 T_N^{*}x=\sum_{j=0}^Np_j\frac{b_j^*x}{\omega_j},\qquad
 T_NT_N^{*}=K_N,
 \tag{CV.6}
\]
where the star uses the fixed Euclidean pairing on the written target
coordinates, solely to specify this typed adjoint.
The canonical right inverse is $T_N^{*}G_N$, whose original source
Gram is $G_N$. Taking its $q$th exterior power maps the fixed
unscaled coordinate wedge to a vector of squared norm
$\det G_N=D_N^{-1}$ in $\bigwedge^q\mathcal P_N$.
The exact exterior source map is
\[
 \bigwedge^q\mathcal V_{h,k}:\bigwedge^q\mathcal P_N
 \longrightarrow\bigwedge^q(\mathscr B^{\widehat\otimes k}).
\]
It preserves the determinant norm $D_N^{-1}$, because the original
source map preserves every pairwise integral inner product.
To pass to the original tensor cochain source, use instead
$\mathcal V_{h,k}^{\otimes q}\operatorname{Alt}_q$, where
$\operatorname{Alt}_q(v_1\wedge\cdots\wedge v_q)
=\sum_{\pi\in S_q}\operatorname{sgn}(\pi)
v_{\pi(1)}\otimes\cdots\otimes v_{\pi(q)}$.
The exterior quotient satisfies
$\operatorname{wed}\operatorname{Alt}_q=q!I$: every permuted
summand contributes its permutation sign twice. Expanding the tensor
inner product gives $q!$ copies of its determinant, so the literal
tensor-source squared norm is $q!D_N^{-1}$. Its target is
$\mathscr B^{\widehat\otimes kq}$ in top cochain degree $kq$.
This proves the asserted interpolation-volume
interpretation without changing its target metric or source mass.
Replacing $T_N$ by its restriction to the omitted-direction subspace
gives exactly $\widetilde K_N$; the construction needs no assertion
that this restricted map is onto.

Now take the packet polynomial to be exactly
\[
 h(s)=\prod_{\varepsilon,\eta\in\{+1,-1\}}
 \bigl(s-(\tfrac12+\varepsilon\delta+i\eta\gamma)\bigr)^m
\]
for a complete actual nonreal quartet
$\frac12\pm\delta\pm i\gamma$, $\delta,\gamma>0$, of common order
$m$, and set $\ell_k=1+k(m-1)$. Its cyclic annihilator has
$q=\ell_k(k+1)^2$. The complete exterior trace calculation gives
$L_k\leq\epsilon$, where
\[
 L_k=2\delta\ell_k(k+1)
              \left\lfloor\frac{(k+1)^2}{4}\right\rfloor.
 \tag{CV.7}
\]
Substitution in (CV.4), with all quantities nonnegative, proves the
exact necessary arithmetic-volume inequality
\[
 \boxed{
 D_{N+1}+D_{N-1}-D_N-\widetilde D_N
 \ \geq\ \frac{\omega_N}{\omega_{N+1}}D_N L_k^2,
 \qquad N\geq\ell_k(k+1)^2-1.}
 \tag{CV.8}
\]
The cross term has been transferred to the actual omitted-direction
determinant by (CV.3); it has not been discarded. Thus (CV.4) supplies
an exact expression for the still unestimated quantity
\[
 \frac{\epsilon^2}{k^6}
 =\frac{\omega_{N+1}}{k^6\omega_N}
 \frac{D_{N+1}+D_{N-1}-D_N-\widetilde D_N}{D_N}.
 \tag{CV.9}
\]
Every term comes from the specified arithmetic convolution measure and
the complete quotient. No decay rate for (CV.9) is asserted here.
The formula identifies which signed determinant combination, and which
orthogonal-norm ratio at growing degree, the next analytic calculation
has to estimate together.
For a larger reflection-stable packet containing this quartet, the
quartet sums remain present and their full nilpotent lengths are at
least $\ell_k$. Indeed the same ordered quartet factors embed as
CRT summands of the larger tensor algebra, and the largest nilpotent
length at a fixed sum is a maximum over all such summands. Their
positive trace contribution is therefore at least (CV.7).
Inequality (CV.8) remains valid, with every $K_j,D_j,\omega_j$
formed from that larger packet and the admission threshold replaced
by its actual $N\geq\deg\chi_{h,k}-1$. No equality for that larger
degree is inferred from the quartet subpacket.

\subsection{The positive exterior minor sum and its actual source vector}

Write
$\Delta_N=D_{N+1}+D_{N-1}-D_N-\widetilde D_N$.
For an increasing finite index set $J$, let $B_J$ be the matrix
whose columns are the original $b_j$, in that increasing order, and
put $\Omega_J=\prod_{j\in J}\omega_j$, with the empty product equal
to one. For every $q\geq2$ and $N\geq q-1$,
\[
 \boxed{
 \Delta_N=
 \sum_{\substack{I\subseteq\{0,\ldots,N-1\}\\ |I|=q-2}}
 \frac{\left|\det B_{I\cup\{N,N+1\}}\right|^2}
      {\omega_N\omega_{N+1}\prod_{i\in I}\omega_i}.}
 \tag{CV.10}
\]
For $q=2$ the index set $I$ is empty and the sum has its single
displayed two-column term. For $q=1$ set this sum to zero; then
$\Delta_N=0$ exactly. Thus the signed combination in (CV.4) is an
explicit sum of positive exterior contributions, retaining every
original column and squared polynomial norm.

\begin{proof}
First take any finite set $L$ of column indices and form
$K_L=\sum_{j\in L}b_jb_j^*/\omega_j$. Multilinearity of its
$q$ columns expands its determinant as
\[
 \det K_L=
 \sum_{(j_1,\ldots,j_q)\in L^q}
 \det[b_{j_1},\ldots,b_{j_q}]
 \prod_{r=1}^q\frac{\overline{(b_{j_r})_r}}{\omega_{j_r}}.
\]
If two indices coincide, the determinant in that summand is zero.
For a fixed increasing set $J=\{i_1<\cdots<i_q\}$, the remaining
terms are indexed by its permutations. Their sum is
\[
 \frac{\det B_J}{\Omega_J}
 \sum_{\sigma\in\mathfrak S_q}\operatorname{sgn}(\sigma)
       \prod_{r=1}^q\overline{(b_{i_{\sigma(r)}})_r}
 =\frac{|\det B_J|^2}{\Omega_J}.
\]
This proves directly, including singular $K_L$,
$\det K_L=\sum_{J\subseteq L,\ |J|=q}|\det B_J|^2/\Omega_J$.
Apply it to the four index sets
\[
 \{0,\ldots,N+1\},\quad \{0,\ldots,N-1\},\quad
 \{0,\ldots,N\},\quad \{0,\ldots,N-1,N+1\}.
\]
For each $q$-element subset of the first set, the coefficients in
the signed combination defining $\Delta_N$ are respectively
$1+1-1-1=0$ if it contains neither $N$ nor $N+1$,
$1-1=0$ if it contains exactly one of them, and $1$ if it
contains both. The surviving subsets have exactly the form
$I\cup\{N,N+1\}$ in (CV.10). If $q=1$, no subset contains
both indices and the result is zero. No inverse for a smaller
or omitted-direction Gram matrix has occurred, so this proof
also applies at $N=q-1$.
\end{proof}

There is a precise source map for the positive sum. Let
$e_C=[1]\wedge[S]\wedge\cdots\wedge[S^{q-1}]$ be the fixed
coordinate wedge in $\bigwedge^qC$, of squared norm one for the
coordinate Euclidean pairing used in (CV.6). For $q\geq2$ define
\[
 \mathcal E_N=(\bigwedge^{q-2}\mathcal P_{N-1})
                      \wedge p_N\wedge p_{N+1}
                 \subseteq\bigwedge^q\mathcal P_{N+1},
 \qquad
 \Lambda_N=(\bigwedge^qT_{N+1})|_{\mathcal E_N}.
\]
The increasing vectors
$e_I=(\bigwedge_{i\in I}p_i)\wedge p_N\wedge p_{N+1}$
form an orthogonal basis of $\mathcal E_N$ with their original
squared norms $\Omega_{I\cup\{N,N+1\}}$. They satisfy
$\Lambda_N e_I=\det B_{I\cup\{N,N+1\}}\,e_C$.
Consequently the adjoint in these explicitly specified source
and target pairings gives the original, unscaled vector
\[
 \begin{split}
 \zeta_N:=\Lambda_N^*e_C
 &=\sum_{\substack{I\subseteq\{0,\ldots,N-1\}\\ |I|=q-2}}
 \frac{\overline{\det B_{I\cup\{N,N+1\}}}}
      {\Omega_{I\cup\{N,N+1\}}}\,e_I,\\
 \|\zeta_N\|^2&=\Delta_N,\qquad
 \Lambda_N\zeta_N=\Delta_N e_C,\\
 \|\mathcal V_{h,k}^{\otimes q}
                \operatorname{Alt}_q\zeta_N\|^2
 &=q!\Delta_N.
 \end{split}
 \tag{CV.10a}
\]
The first two equalities follow by inserting the displayed
orthogonal basis and its full norms. The last equality follows
from the pairwise isometry of $\mathcal V_{h,k}$ and the
unscaled alternating-map calculation following (CV.6). Its
target is the original tensor source
$\mathscr B^{\widehat\otimes kq}$ in degree $kq$.
For $q=1$ take $\mathcal E_N=\{0\}$, $\Lambda_N=0$, and
$\zeta_N=0$, so that all three zero identities remain exact.
Thus (CV.10) comes with an explicitly typed original arithmetic
source vector and the literal tensor factorial. Combining it
with (CV.4) gives
$\epsilon^2=(\omega_{N+1}/\omega_N)\|\zeta_N\|^2/D_N$;
no estimate as $k$ or $N$ grows is inferred from this equality.

\subsection{The exact first admitted degree}

At $N=q-1$ expand the original monic annihilator in the original
orthogonal basis:
\[
 \chi=p_q+\sum_{j=0}^{q-1}\gamma_jp_j,\qquad
 \gamma_j=\frac{\int\overline{p_j(k/2+iu)}\chi(k/2+iu)m_{h,k}(u)\,du}{\omega_j}.
 \tag{CV.11}
\]
This is a finite polynomial identity: both leading coefficients are one,
and orthogonality gives every displayed coefficient, without a change of
measure. The reduction map $T_{q-1}$ is a linear isomorphism, since the
$q$ original monic remainders are a basis. Its inverse is therefore its
canonical minimum section. Consequently
$b_i^*G_{q-1}b_j=\delta_{ij}\omega_j$ for $0\le i,j<q$.
Reduction of (CV.11) gives $b_q=-\sum_{j<q}\gamma_jb_j$ and hence
\[
 \begin{gathered}
 a=\omega_{q-1},\quad
 c=-\gamma_{q-1}\omega_{q-1},\quad
 d=\sum_{j<q}|\gamma_j|^2\omega_j,\\
 \boxed{\epsilon_{h,k,q-1}^2
 =\frac{\sum_{j=0}^{q-2}|\gamma_j|^2\omega_j}{\omega_{q-1}}
 =\frac{\|\chi\|^2-\omega_q-|\gamma_{q-1}|^2\omega_{q-1}}{\omega_{q-1}}.}
\end{gathered}
 \tag{CV.12}
\]
The first expression follows by substituting the preceding three exact
Gram entries into (CV.4); the second follows by Pythagoras applied to
(CV.11). Reflection makes $c$ purely imaginary, so the same is true of
$\gamma_{q-1}$; it has not been set to zero. Let $\mathsf P_{q-2}$
be the original integral orthogonal projection onto
$\mathcal P_{q-2}$, with $\mathcal P_{-1}=0$ for $q=1$.
Then the numerator in (CV.12) is exactly
$\|\mathcal V_{h,k}\mathsf P_{q-2}\chi\|^2$ by the original
source isometry. Meanwhile the full $\mathcal V_{h,k}\chi$ has
zero cyclic and full jets, since $\chi(S)\in I$, and the original
fixed-order polynomial division gives its theta primitive as in
(CM.19). These are different specified source maps: projecting that
boundary is an integral orthogonal operation and has the finite
jet $\sum_{j<q-1}\gamma_jb_j$. This exact formula retains its
possibly nonzero observation. For $q=1$ the sum is empty and both
sides are zero. The degree-$q$ annihilator and its projection have
thus given a complete original-source formula at the first admitted
degree, with no large-degree estimate added.
